\section{Repaso de las diapositivas clave}

\subsection{Portada y posicionamiento del tema}

\begin{figure}[H]
\centering
\includegraphics[width=0.8\textwidth]{cover.jpg}
\caption{Portada del curso: revisión de las tendencias de GenAI desde la perspectiva empresarial}
\end{figure}

La portada señala: esta clase no es una investigación de algoritmos, sino que aborda el tema desde cuatro niveles empresariales: GTM, plataforma, herramientas y gobernanza. Esto también determina que, en todo el conjunto de diapositivas, cada tema incluya cuatro dimensiones: casos, datos, cadena de herramientas y políticas.

\subsection{Mapa de tendencias y timeline}

En varias partes de las diapositivas aparece el mapa de tendencias: en el lado izquierdo se enumeran las capability técnicas (como multi-modal, agent), y en el lado derecho las necesidades empresariales (como compliance, observabilidad). Burak insiste repetidamente en colgar esta imagen en la pared, para que el equipo la consulte al elaborar el roadmap.

\begin{knowledgebox}{Cómo usar el mapa de tendencias}
\begin{enumerate}
    \item Primero ubica tu capability: ¿cuentas con un search + tool loop?
    \item Después revisa el lado derecho de las necesidades: ¿qué regulator, qué SLA aún no se han cumplido?
    \item Por último, dibuja el gap: el gap en sí mismo es la base de la prioritization
\end{enumerate}
\end{knowledgebox}

\subsection{Demostración de tiempo real y observabilidad}

Otra diapositiva clave muestra el dashboard de observabilidad de Vertex Agent: en el lado izquierdo está latency/time series, y en el lado derecho está tool-call matrix. Burak señala que el dashboard no es un nice-to-have, sino una de las condiciones del product launch.

\subsection{Resumen de la sección}

Las diapositivas no solo ofrecen un marco macro, sino que también pueden reutilizarse directamente en el workshop del equipo (mapa de tendencias + plantilla de dashboard).

